\subsubsection{Two Micron-Size Dark Dimensions --- arXiv:2501.11690}

\textbf{Authors:} Luis A. Anchordoqui, Ignatios Antoniadis, Dieter L\"ust

\paragraph{主な主張}
ミクロンサイズの1次元・2次元 large extra dimensions に対する collider・天体物理・宇宙論的制限を系統的に検討し、2次元の場合は astrophysical constraints を満たすために isometry を持たず KK momentum conservation を破って重い KK graviton の lighter modes への崩壊を可能にする必要があり、さらに cosmological constraints との両立には radiation-dominated epoch の開始温度の微調整が必要であることを示した \cite{Anchordoqui:2025TwoMicron}。

\paragraph{新規性}
two micron-size dark dimensions の phenomenological viability を複数の観測制限から横断的に評価し、isometry の欠如による KK graviton の intra-tower decay を天体物理制限を回避するための中心的条件として明確化した。

\paragraph{位置付け}
ADD / Dark Dimension の2次元 compactification が成立するための全体的な consistency conditions を整理した研究であり、$T^2$ など具体的な2次元幾何で KK spectrum や neutrino phenomenology を議論する際の前提条件を与える。
